\phantomsection
\chapter*{Kampung Halaman}
\addcontentsline{toc}{section}{Kampung Halaman --- Hendrik Ramadani}
\penulis{Hendrik Ramadani}


\awalcerpen{S}{udah enam tahun} saya tidak menginjakkan kaki ke gang sempit yang penuh kenangan ini. Jalan yang dulu retak-retak sekarang sudah diperbaiki. Pohon mangga yang dulu kecil sekarang sudah menjadi besar dan berbuah; di tempat yang sama, aku dan Ihsan setiap pulang mengaji selalu memanjat pohon mangga itu setiap sore. 

Setelah mengingat kenangan dulu, aku terhenti sejenak di depan rumah Ihsan. Catnya sudah berubah dari biru ke putih. Aku menarik napas dalam-dalam untuk bersiap bertemu Ihsan, lalu aku mengetuk pintu. Tak kusangka, yang membuka pintu adalah seorang nenek yang tersenyum ramah. Namun, ia tidak mengenali aku. Nenek itu bertanya, ``Cari siapa ya, Mas?'' 

Dan aku menjawab, ``Aku pernah tinggal di sini, Nek. Saya sedang mencari seseorang yang bernama Ihsan yang dulu rumahnya di sini.''

Nenek itu tampak berpikir sejenak untuk mengingat nama Ihsan. Nenek pun menjawab, ``Oh, Ihsan anaknya Pak Karim? Ihsan dan keluarganya sudah pindah ke Jagabaya lama sekali. Tapi kalau Mas mau mencari rumah Ihsan, nenek bisa bantu arahkan. Dari sini lurus terus, belok kiri sampai ketemu masjid, lalu belok kanan. Cari saja rumah yang catnya berwarna hijau.''

Aku mengangguk pelan dan berterima kasih atas arahan yang diberikan oleh nenek. Lalu, aku menyusuri gang yang dulu sangat aku hafal. Aku ingat betul setiap pulang sekolah aku dan Ihsan selalu mampir ke rumah Junior untuk membeli \textit{pop ice} yang harganya tiga ribu rupiah, lalu bermain kelereng dan galah. Ihsan jago banget bermain kelereng.

Aku menyusuri jalan sesuai arahan dari nenek untuk mencari rumah Ihsan yang berwarna hijau. Di tengah jalan, karena iseng, aku memutuskan pergi ke rumah Junior terlebih dahulu untuk membeli \textit{pop ice} seharga tiga ribu rupiah. Tak menyangka, warung Ibu Junior masih berdiri, meski sudah direnovasi menjadi sangat rapi. Aku masuk ke warung Ibu Junior untuk memesan \textit{pop ice}. 

Sambil menunggu pesanan \textit{pop ice}, ponselku tiba-tiba bergetar dari \textit{number} yang tak dikenal. Setelah aku angkat, suara di seberang tiba-tiba memanggil namaku, ``Karlan?''

Aku sempat hendak menutup telepon, namun tiba-tiba suara itu melanjutkan, ``Ini aku, Ihsan. Ibu Junior cerita kepadaku kalau ada yang mencari kamu, dan ciri-cirinya persis seperti teman masa kecilku.'' 

Ihsan kemudian memanggilku lagi di telepon, ``Karlan, masih di warung Ibu Junior nggak? Aku \(\textit{otw}\) kesana, jangan pulang dulu!''

Aku tersenyum sendiri menatap layar ponsel. Di seberang sana, Ihsan langsung berlari menuju rumah Junior untuk menemuiku.

\jedahias
